\documentclass[11pt,a4paper]{article}

\usepackage[T1]{fontenc}
\usepackage[utf8]{inputenc}
\usepackage[spanish]{babel}
\usepackage{geometry}
\usepackage{microtype}
\usepackage{setspace}
\usepackage{booktabs}
\usepackage{tabularx}
\usepackage{enumitem}
\usepackage{hyperref}
\usepackage{xcolor}
\usepackage{titlesec}
\usepackage{fancyhdr}
\usepackage{array}
\usepackage{parskip}
\usepackage{tikz}

\usetikzlibrary{arrows.meta,positioning}

\geometry{margin=2.3cm, top=2.6cm, bottom=2.6cm}
\setstretch{1.15}
\setlist[itemize]{leftmargin=1.4em,itemsep=0.28em,topsep=0.4em}
\setlist[enumerate]{leftmargin=1.4em,itemsep=0.28em,topsep=0.4em}

\definecolor{upnavy}{HTML}{1F3A5F}
\definecolor{upgray}{HTML}{5A6570}
\definecolor{upsoft}{HTML}{E8EEF4}

\hypersetup{colorlinks=true,urlcolor=upnavy,linkcolor=upnavy}
\titleformat{\section}{\large\bfseries\color{upnavy}}{\thesection.}{0.6em}{}
\titleformat{\subsection}{\normalsize\bfseries\color{upnavy}}{\thesubsection.}{0.5em}{}
\titleformat{\subsubsection}{\normalsize\itshape\color{upnavy}}{\thesubsubsection.}{0.5em}{}

\pagestyle{fancy}
\fancyhf{}
\fancyhead[L]{\small\color{upgray}Pipeline Drive --- conflictos Defensoría}
\fancyhead[R]{\small\color{upgray}Universidad del Pacífico}
\fancyfoot[C]{\thepage}
\renewcommand{\headrulewidth}{0.3pt}

\newcolumntype{Y}{>{\raggedright\arraybackslash}X}

\begin{document}

\begin{titlepage}
\centering
\vspace*{1.4cm}
{\color{upnavy}\rule{\textwidth}{1.4pt}}\\[0.7em]
{\LARGE\bfseries\color{upnavy} El pipeline de Drive:}\\[0.35em]
{\LARGE\bfseries\color{upnavy} de PDF a 83 conflictos mineros}\\[0.7em]
{\Large Carpeta Minería --- disco \texttt{Johao\_asistente}\\[0.2em]
(2024 -- febrero 2025)}\\[0.7em]
{\color{upnavy}\rule{\textwidth}{1.4pt}}\\[1.3em]
{\large Cómo se construyó el CSV que el grupo usaba\\[0.2em]
y por qué ese flujo se dejó de lado}\\[2em]

\begin{flushleft}
\renewcommand{\arraystretch}{1.35}
\begin{tabular}{@{}l@{\hspace{1.2em}}l@{}}
\textbf{Proyecto} & Minería, conflictos sociales e indicadores territoriales \\
\textbf{Institución} & Universidad del Pacífico \\
\textbf{Fuente} & Reportes mensuales de conflictos sociales, Defensoría del Pueblo \\
\textbf{Disco Drive} & \texttt{Johao\_asistente} $\rightarrow$ carpeta Minería \\
\textbf{Código} & \texttt{02\_Codigo/conflictos\_defensoria/Extraccion\_2004\_2025} \\
\textbf{Salidas} & \texttt{Conflicto.csv} (83 casos); \texttt{ReporteFechaxConflicto.csv} (3\,777 filas) \\
\textbf{Cobertura} & 251 ediciones (n.$^\circ$~1--251, hasta enero 2025) \\
\textbf{Fecha de esta nota} & Septiembre 2026 \\
\end{tabular}
\end{flushleft}

\vfill
\begin{flushleft}
\color{upgray}\small
Este documento describe \emph{solo} el pipeline viejo de Drive: pasos, herramientas y puntos en contra, con evidencia de los notebooks y del Word de anotaciones del propio equipo. No es un paper y no se compiló a propósito; vive en \texttt{04\_docs}. El contraste detallado con el extractor local (\texttt{defensoria\_extract}) está en \texttt{metodologia\_extraccion\_conflictos.tex}.
\end{flushleft}
\end{titlepage}

\tableofcontents
\newpage

%----------------------------------------------------------------------
\section{En una frase}

El CSV de la carpeta Minería \textbf{no es el universo de conflictos de la Defensoría}. Es el resultado de recortar PDF por rúbricas del índice, convertirlos a Excel con iLovePDF, parsear cada ventana 2004--2025 con código distinto, agrupar textos con fuzzy nacional y quedarse con \textbf{83 casos} que coincidían con un catálogo cerrado de minas (más revisión manual y un resumen de ChatGPT).

Eso es un producto de \emph{investigación} (universo minero del grupo). No es un volcado exhaustivo de fichas. Quien lo trate como ``la base de conflictos sociales'' está leyendo mal el método.

%----------------------------------------------------------------------
\section{Mapa del flujo}

\begin{center}
\begin{tikzpicture}[
  font=\footnotesize,
  box/.style={draw=upnavy, rounded corners=2pt, align=center, inner sep=5pt,
              fill=upsoft, text=upnavy, minimum height=0.95cm, text width=1.85cm},
  arr/.style={-{Latex}, thick, color=upnavy}
]
\node[box] (a) {PDF\\n.$^\circ$~1--251};
\node[box, right=0.35cm of a] (b) {Recorte\\por frase};
\node[box, right=0.35cm of b] (c) {iLovePDF\\$\to$ Excel};
\node[box, right=0.35cm of c] (d) {Parser\\por periodo};
\node[box, below=0.7cm of a] (e) {Unificación\\de columnas};
\node[box, below=0.7cm of b] (f) {Fuzzy\\nacional 90};
\node[box, below=0.7cm of c] (g) {Catálogo\\minas + mano};
\node[box, below=0.7cm of d] (h) {83 casos\\+ panel 3\,777};
\draw[arr] (a) -- (b);
\draw[arr] (b) -- (c);
\draw[arr] (c) -- (d);
\draw[arr] (d) |- (e);
\draw[arr] (e) -- (f);
\draw[arr] (f) -- (g);
\draw[arr] (g) -- (h);
\end{tikzpicture}
\end{center}

Ruta en Drive: \texttt{Minería/02\_Codigo/conflictos\_defensoria/Extraccion\_2004\_2025}. Notebooks centrales: \texttt{Procesamiento 2004 - 2025.ipynb} ($\approx$498 celdas), \texttt{Unificación.ipynb}, \texttt{Fuzzy.ipynb}, \texttt{Casos\_Filtrados\_Actualizado.ipynb}. Inventario: \texttt{conflictos\_sociales\_pdfs.csv} (251 ediciones, todas marcadas \texttt{Revisado=OK}).

%----------------------------------------------------------------------
\section{Cómo funcionaba}

\subsection{Corpus e inventario}

Había 251 PDF mensuales (abril 2004 -- enero 2025). El estado del conflicto ---Activo, Latente, Resuelto, Inactivo, Fusionado, Nuevo--- \textbf{no se leía de la ficha}. Se heredaba de la carpeta a la que caía el recorte: una sección del índice = un archivo = un estado impuesto al concatenar.

\subsection{Paso 0. Conversión externa (iLovePDF)}

El código no leía el PDF como documento con geometría. El flujo, visible en rutas del estilo \texttt{C:\textbackslash Users\textbackslash Rofer\textbackslash \ldots\textbackslash ILOVEPDF\textbackslash \ldots}, era:

\begin{enumerate}
  \item Recortar el PDF original a un PDF ``de sección'' (solo activos, solo latentes, etc.).
  \item Pasar ese recorte por iLovePDF (o equivalente) a Excel o Word.
  \item Parsear el Excel/Word con pandas o \texttt{python-docx}.
\end{enumerate}

Consecuencia de método: \textbf{se pierde la geometría de página}. Las dos columnas de una ficha (descripción a la izquierda, hechos del mes a la derecha) llegan como celdas de una hoja, no como bloques con coordenadas $(x,y)$. Si el conversor mezcla columnas, el parser hereda el error.

Herramientas de recorte: \texttt{pdfplumber} para buscar frases de inicio/fin; \texttt{PyPDF2} (\texttt{PdfReader}/\texttt{PdfWriter}) para copiar el rango a \texttt{recortado\_*.pdf}.

\subsection{Paso 1. Recorte por rúbricas del índice}

La función central era \texttt{recortar\_pdfs(frase\_inicio, frase\_fin, \ldots)}. Recorría páginas, guardaba el índice de la primera que contenía la frase de inicio y el de la primera posterior con la de fin, y escribía un PDF nuevo. Las frases se reescribían a mano cada vez que la Defensoría renumeraba el índice. Ejemplos del notebook 2024--2025:

\begin{itemize}
  \item Activos, setiembre 2024: ``V.\ DETALLE DE LOS CONFLICTOS SOCIALES ACTIVOS'', luego ``5.1 \ldots EN UN SOLO DEPARTAMENTO'', corte en ``5.3 CONFLICTOS REACTIVADOS''.
  \item Latentes: ``VI.\ DETALLE DE LOS CONFLICTOS LATENTES'' / ``6.1 \ldots DE ESTADO ACTIVO A LATENTE''.
  \item Inactivos: ``6.2 \ldots POR INACTIVIDAD PROLONGADA''.
  \item Resueltos: ``1.4 CONFLICTOS RESUELTOS'' (en octubre 2024 ya era ``8.1'' dentro de ``VIII.\ CASOS QUE SALIERON DEL REPORTE'').
\end{itemize}

Hubo once periodos con parser propio porque la maqueta no era estable:

\begin{center}
\renewcommand{\arraystretch}{1.15}
\begin{tabularx}{\textwidth}{@{}l Y@{}}
\toprule
\textbf{Periodo} & \textbf{Qué se recortaba y cómo se leía} \\
\midrule
Abril 2004 & Solo activos; Word (\texttt{python-docx}); dos casos de Cajamarca a mano. \\
Mayo 2004 & Excel de tabla regional; departamento al ver \texttt{REGIÓN \ldots}. \\
Jun--ago 2004 & Activos, nuevos y concluidos; Antecedentes / Cuestionamiento / Hechos. \\
Set--oct 2004 & Se suman inactivos (``Desactivado'' vs ``Resuelto''). \\
Nov 2004--dic 2005 & Activos, latentes y resueltos. \\
Ene--dic 2006 & Texto corrido; funciones distintas por estado. \\
Ene 2007--mar 2008 & Tablas + texto (reportes 35--49). \\
Abr 2008--jul 2011 & Cuatro estados; \texttt{pdfplumber} sistemático. \\
Ago 2011--jul 2014 & Latentes en dos oleadas (hasta n.$^\circ$~108 y desde 109). \\
Ago--set 2014 & Se añaden fusionados. \\
Oct 2014--set 2024 & Maqueta Adjuntía; parser de dos columnas sobre Excel. \\
Oct 2024--ene 2025 & Mismos campos; rúbricas de índice distintas (ciclo). \\
\bottomrule
\end{tabularx}
\end{center}

\subsection{Paso 2. Parseo del Excel o Word}

Una vez el recorte era una hoja, el parser típico de Adjuntía (\texttt{procesar\_reporte\_conflictos\_activos}) hacía esto:

\begin{enumerate}
  \item Leía \emph{todas} las hojas (\texttt{sheet\_name=None}, sin encabezado).
  \item Tiraba filas cuyo texto coincidía con títulos de sección o con ``Descripción'', ``Hechos del mes'', ``Estado actual''.
  \item Forzaba dos columnas: izquierda = ficha, derecha = hechos (a veces ``todo menos la última'', a veces primera y última).
  \item Acumulaba filas hasta ver un nuevo \texttt{Tipo} / \texttt{CASO} / departamento; entonces cerraba la ficha anterior.
  \item Partía la descripción con regex (\texttt{procesar\_parrafo}); si hacía falta, \texttt{combinar\_caso\_tipo}.
  \item Detectaba diálogo en la columna de hechos (\texttt{procesar\_columna\_dialogo}).
\end{enumerate}

Inicio de ficha: celda izquierda que empieza por \texttt{Tipo} o \texttt{CASO}, o un departamento de una lista de 25 nombres. En 2004--2007 el criterio era otro: filas numeradas (\texttt{1.\ \ldots}), bloques \texttt{REGIÓN X}, y el campo \texttt{Caso} a menudo se \emph{reconstruía} concatenando Antecedentes + Hechos + Cuestionamiento.

\subsection{Paso 3. Unificación de esquemas}

\texttt{Unificación.ipynb} no volvía a leer PDF. Cargaba un Excel por periodo y por estado, renombraba columnas a un esquema común y concatenaba: \texttt{activototal}, \texttt{latentetotal}, \texttt{inactivototal}, \texttt{resueltototal}, \texttt{fusionadototal}. El campo \texttt{Estado} se imponía al concatenar (\texttt{activototal['Estado']='Activo'}), no se leía del PDF.

Ejemplos de armonización (celda a celda, no un diccionario único): mayo 2004, \texttt{Motivo} $\to$ Caso; set--oct 2004, Caso = Antecedentes + Hechos + Cuestionamiento; latentes 2008--2011, tipo separado con una lista de 30+ variantes.

Salida intermedia: \texttt{ConflictosTotales2024\_2025.xlsx} (todos los estados, todavía sin identidad fuzzy).

\subsection{Paso 4. Identidad fuzzy}

\texttt{Fuzzy.ipynb}:

\begin{enumerate}
  \item Estandarizaba \texttt{Caso} a minúsculas, quitaba puntuación, \textbf{conservaba tildes} (NFC), no quitaba stopwords.
  \item Cruzaba con el inventario de PDF para pegar fecha y número de edición.
  \item Agrupaba con \texttt{rapidfuzz.fuzz.token\_sort\_ratio}, umbral \textbf{90}, sobre la \textbf{lista nacional} de textos (complejidad $n^2$, sin bloquear por departamento).
  \item Asignaba \texttt{Grupo\_k}; calculaba primera/última fecha, meses faltantes y tránsitos de estado.
\end{enumerate}

Límite documentado por el propio equipo: Excel trunca celdas en 32\,767 caracteres. Nota en \texttt{Anotaciones - Detalles a mejorar.docx}: ``Hay una limitación de 32767, para agrupar''.

\subsection{Paso 5. Filtro de minería (lista cerrada + revisión)}

\texttt{Casos\_Filtrados\_Actualizado.ipynb} no usaba el tipo ``socioambiental'' como filtro principal. Limpiaba tildes con \texttt{unidecode} y buscaba \textbf{palabras exactas} de una lista armada a mano:

\begin{itemize}
  \item Unidades / compañías: Antamina, Antapaccay, Atacocha, Cerro Corona, Las Bambas, Pucamarca, Quellaveco, San Cristóbal, Toromocho, Yanacocha, Glencore, Nexa, Gold Fields, MMG, Minsur, Anglo, Volcan, Chinalco, Newmont, Xstrata.
  \item Pueblos asociados (San Marcos, Challhuahuacho, Hualgayoc, Nueva Morococha, \ldots).
\end{itemize}

Quien matcheaba pasaba a \texttt{FuzzyCompletar\_filtrado}. Después: etiquetado manual (\texttt{EsMinera}, \texttt{EsIlegal}, \texttt{EsInformal}, pueblos, minas, compañía) y resumen ChatGPT. El CSV final \texttt{Conflicto.csv} tiene 83 filas, \textbf{todas} \texttt{EsMinera = Sí}.

El panel \texttt{ReporteFechaxConflicto.csv} no es el universo de la Defensoría: es la serie de esos 83 grupos (3\,777 filas; 3\,208 activos, 508 latentes, 34 resueltos, 26 inactivos, 1 fusionado; 244 PDF distintos; ediciones 7--251).

\subsection{Caja de herramientas}

iLovePDF (fuera del repo), \texttt{pdfplumber}, \texttt{PyPDF2}, pandas, \texttt{python-docx}, \texttt{rapidfuzz}, \texttt{unidecode}, ChatGPT, notebooks ($\approx$498 celdas en el principal), rutas de un PC concreto (\texttt{C:\textbackslash Users\textbackslash Rofer\textbackslash \ldots}).

%----------------------------------------------------------------------
\section{Puntos en contra}

Cada punto: qué se hacía, por qué duele, y un ejemplo concreto (del Word de anotaciones o de los números del propio producto).

\subsection{1. El conversor aplana el layout}

\textbf{Qué.} iLovePDF convierte el PDF de sección a celdas; el parser adivina cuál columna es hechos del mes.

\textbf{Por qué duele.} Una familia entera de errores nace ahí: tabulaciones que se comen latentes, el match de ``Actor'' que corta la palabra ``directorio'', secciones que no agarran el párrafo entero. El parser no puede saber si el conversor mezcló izquierda y derecha.

\textbf{Ejemplo.} Latentes marzo 2006--marzo 2007 incompletos por error de tabulación; latentes 2007--2008 incompletos porque la sección no agarra el párrafo.

\subsection{2. Once tijeras en vez de un detector}

\textbf{Qué.} Frases literales del índice, reescritas por periodo (``V.\ DETALLE\ldots'' vs ``VI.\ DETALLE\ldots''; ``1.4'' vs ``8.1'').

\textbf{Por qué duele.} Si la Defensoría renumeraba el capítulo, el recorte fallaba en silencio o había que abrir el notebook otra vez. Costo operativo: $\approx$498 celdas y once parsers.

\textbf{Ejemplo.} Resueltos enero 2007--marzo 2008 se recortaron con la frase de latentes (anotación del propio equipo).

\subsection{3. El estado es de la carpeta, no de la ficha}

\textbf{Qué.} Activo / Latente / Resuelto se impone al concatenar según el bloque (\texttt{activototal}, \texttt{latentetotal}, \ldots).

\textbf{Por qué duele.} Un recorte mal etiquetado etiqueta mal \emph{todas} las filas de ese bloque. No hay forma de auditar el estado mirando solo el texto de la ficha.

\textbf{Ejemplo.} El mismo episodio de resueltos recortados como latentes: el error de frase se convierte en error de variable de análisis.

\subsection{4. La ficha no cruza bien de página}

\textbf{Qué.} El recorte parte el PDF en rangos de sección. Una ficha que empieza al final de un rango y sigue en el siguiente se pierde o se parte.

\textbf{Por qué duele.} El parser de Excel acumula celdas hasta ver otro \texttt{Tipo} o un departamento; si el conversor o el recorte cortan a mitad, se generan filas rotas.

\textbf{Ejemplo.} Faltan activos de mayo 2006 (n.$^\circ$~27) ``por el inicio de la página inicial'' (Word de anotaciones).

\subsection{5. Identidad nacional, umbral 90}

\textbf{Qué.} \texttt{token\_sort\_ratio} $\ge$ 90 sobre la lista nacional; tildes conservadas; sin stopwords; sin bloquear por departamento.

\textbf{Por qué duele.} Mezcla homónimos de distintos departamentos. Complejidad $n^2$. Conservar tildes y no quitar stopwords hace el match más frágil cuando la Defensoría reescribe el mismo caso mes a mes (``el caso de la comunidad \ldots'' vs ``comunidad \ldots'').

\textbf{Ejemplo.} Corire no debería compartir caso con Racrachaca (anotado por el equipo). Diciembre 2012 ``no corrió bien''.

\subsection{6. Minería = lista cerrada + borrado del resto}

\textbf{Qué.} Filtro por nombres de unidades, compañías y pueblos; luego etiqueta manual \texttt{EsMinera=Sí}. Quedan 83 filas.

\textbf{Por qué duele.} Un conflicto socioambiental artesanal, o una unidad que no está en el catálogo, no entra ---aunque el tipo diga socioambiental. El CSV no sirve para preguntar ``cuántos conflictos socioambientales hay en el reporte $t$''. Sesga el universo hacia las minas grandes del proyecto (Yanacocha, Antamina, Bambas, Toromocho, Quellaveco, \ldots).

\textbf{Ejemplo.} Tipos residuales (laboral, gobierno local) aparecen en el CSV de 83 porque el filtro no era el tipo Defensoría, sino la lista de minas; lo que no matcheaba la lista simplemente no estaba.

\subsection{7. Huecos conscientes, no auditables en el CSV}

\textbf{Qué.} El Word \texttt{Anotaciones - Detalles a mejorar.docx} lista fallos conocidos y no corregidos al cerrar el producto.

\textbf{Por qué duele.} El CSV se usó igual. Quien abre \texttt{Conflicto.csv} no ve qué meses faltan ni qué secciones se cortaron mal: los huecos viven en un Word aparte.

\textbf{Ejemplos del propio equipo.}
\begin{itemize}
  \item Faltan activos de mayo 2006 (n.$^\circ$~27).
  \item Faltan latentes marzo 2006--marzo 2007.
  \item Resueltos enero 2007--marzo 2008 con frase de latentes.
  \item 2007--2011: ``Actor'' corta ``directorio''.
  \item Latentes 2007--2008 incompletos.
  \item Diciembre 2012 mal; Corire vs.\ Racrachaca.
\end{itemize}

\subsection{8. No reproducible}

\textbf{Qué.} iLovePDF fuera del repositorio; rutas de un usuario; conversión no versionada; resumen ChatGPT de una corrida concreta.

\textbf{Por qué duele.} No se puede volver a correr el pipeline de punta a punta en otra máquina y obtener el mismo CSV. Excel trunca celdas largas (32\,767). El resumen ChatGPT no se puede repetir al pie de la letra.

\textbf{Ejemplo.} Las rutas \texttt{C:\textbackslash Users\textbackslash Rofer\textbackslash \ldots\textbackslash ILOVEPDF\textbackslash \ldots} son parte del método, no un detalle cosmético.

\subsection{9. Cobertura corta y universo chico}

\textbf{Qué.} 251 ediciones (hasta enero 2025); el panel minero empieza en la edición 7; 3\,777 filas caso $\times$ mes solo de 83 grupos.

\textbf{Por qué duele.} Eso no es un bug de parseo: es el filtro. Quien necesita el universo (o ediciones 252--269) no puede usarlo como base. Comparado con el extractor local posterior: 269 ediciones, del orden de 35\,000 filas, minería como bandera y no como borrado.

%----------------------------------------------------------------------
\section{Resumen: decisión, método Drive, costo}

\begin{center}
\renewcommand{\arraystretch}{1.2}
\begin{tabularx}{\textwidth}{@{}l Y Y@{}}
\toprule
\textbf{Decisión} & \textbf{Cómo lo resolvía Drive} & \textbf{Costo} \\
\midrule
Leer el PDF & Recorte + iLovePDF $\to$ Excel/Word & Se pierde $(x,y)$; columnas mezcladas \\
Dónde está el detalle & Frases literales del índice, por periodo & Once tijeras; fallos silenciosos \\
Estado & Carpeta del recorte & Error de frase = error de variable \\
Ficha entre páginas & Frágil (el recorte parte el PDF) & Filas rotas o perdidas \\
Identidad & Fuzzy nacional, umbral 90 & Homónimos; $n^2$; Excel 32\,767 \\
Minería & Lista cerrada + etiqueta manual & 83 casos; se borra el resto \\
Calidad & Word de huecos conocidos & Huecos no visibles en el CSV \\
Reproducir & Rutas de un PC + iLovePDF & No reproducible \\
Cobertura & 251 PDF; panel desde n.$^\circ$~7 & Sin eds.\ 252--269; universo chico \\
\bottomrule
\end{tabularx}
\end{center}

%----------------------------------------------------------------------
\section{Qué sí hacía bien (y no se tira)}

Conviene no vender el abandono del pipeline como si el CSV de 83 no valiera nada.

\begin{itemize}
  \item \textbf{Curación.} Minas, pueblos y compañías etiquetados a mano; resumen legible (\texttt{CasoResumidoChatGPT}). Eso el panel crudo no lo tiene.
  \item \textbf{Campos de tránsito.} Motivo de fusión, pase a latente/inactivo, forma de resolución, diálogo como columna propia: el flujo de Drive los priorizaba.
  \item \textbf{Latentes tabulares.} Intentó extraerlos con recortes específicos; a cambio fallaba cuando la frase no coincidía.
\end{itemize}

El CSV de 83 sigue siendo un producto de \emph{investigación} (universo minero del grupo). El extractor local es el de \emph{extracción} (leer el PDF de forma estable y completa). Se pueden cruzar: filtrar el panel local por mención de minería y, si se quiere, aplicar encima el catálogo de minas de Drive.

%----------------------------------------------------------------------
\section{Cierre}

No se dejó este pipeline por capricho. Se dejó porque el layout, el recorte por frases y el filtro de minas convertían el universo de la Defensoría en \textbf{83 filas irreproducibles}, con huecos que el propio equipo ya había escrito en un Word y que el CSV no mostraba.

La mejora que importa no es un umbral fuzzy distinto. Es dejar de entregar el layout a un conversor, dejar de mantener once tijeras y dejar la decisión ``qué es minería'' para el análisis ---no para el paso que borra el resto.

%----------------------------------------------------------------------
\section{Fuentes}

No se adivinó el método: se leyeron los archivos vigentes en Drive.

\begin{itemize}
  \item \texttt{Procesamiento 2004 - 2025.ipynb} --- recorte, parsers por periodo.
  \item \texttt{Unificación.ipynb} --- rename/concat; escribe \texttt{ConflictosTotales2024\_2025.xlsx}.
  \item \texttt{Fuzzy.ipynb} --- umbral 90, \texttt{FuzzyCompletar.xlsx}.
  \item \texttt{Casos\_Filtrados\_Actualizado.ipynb} --- lista de minas/pueblos.
  \item \texttt{Anotaciones - Detalles a mejorar.docx} --- fallos conocidos.
  \item \texttt{conflictos\_sociales\_pdfs.csv}, \texttt{Conflicto.csv}, \texttt{ReporteFechaxConflicto.csv}.
\end{itemize}

Carpeta Minería: \url{https://drive.google.com/drive/folders/1NfRO7wAfOVF6dasoElMsV8srW-7ThSh9}

Para el contraste con el extractor local: \texttt{04\_docs/metodologia\_extraccion\_conflictos.tex}.

\end{document}
